% !TeX root = main.tex
\newpage
\section{Python Configuration Variables}

\subsection{Purpose And Scope}

This section documents the configuration surface exercised by
\texttt{BoundaryParticleReservoirPeriodic.py}.  The scenario cfg currently
contains settings owned by several different systems:

\begin{itemize}
    \item the data-generation front end;
    \item the particle preview plotter;
    \item the Python simulation game UI;
    \item \texttt{ForceDynamics};
    \item the reservoir particle generator;
    \item generated Vulkan \texttt{.tst} compatibility data;
    \item global dynamics-model selection in \texttt{ParticleUtil.cfg}.
\end{itemize}

The fields are intentionally grouped below by owner.  A field being present in
the cfg does not mean that every subsystem uses it.

\subsection{Configuration Assembly}

The active scenario begins with:

\begin{verbatim}
cfg_gendata/BoundaryParticleReservoirPeriodic.cfg
\end{verbatim}

The optional \texttt{include\_file} is loaded after the parent cfg.  Every key
in the included file replaces the same key from the parent.  Keys absent from
the include retain their parent values.  Therefore the effective configuration
cannot be determined by reading only one of the two files.

The usual selections are:

\begin{verbatim}
include_file = ".../sub_BPRPPy.cfg";  // reduced Python case
include_file = ".../sub_BPRPVk.cfg";  // larger Vulkan case
\end{verbatim}

Only one include assignment should be active.  Active project cfg files use a
flat key layout.  \texttt{get\_run\_configuration(config)} retains temporary
read compatibility with old files containing a \texttt{RUN\_CONFIGURATION}
subgroup.

\subsection{Global Dynamics Selection}

These fields belong to \texttt{ParticleUtil.cfg}, not the scenario cfg:

\begin{description}
    \item[\texttt{live\_base}] Selects the primary Python dynamics model.
    For this scenario the normal value is \texttt{"ForceDynamics"}.

    \item[\texttt{shadow\_base}] Selects an independently running comparison
    model.  Set it to \texttt{"none"} to disable the shadow simulation.

    \item[\texttt{display\_base}] Selects \texttt{"live"} or
    \texttt{"shadow"} for display.  It does not stop either model from
    running.  It must be \texttt{"live"} when \texttt{shadow\_base} is
    \texttt{"none"}.
\end{description}

\subsection{Generator And File Selection}

\begin{description}
    \item[\texttt{type}] Front-end generator category.  The current scenario
    uses \texttt{"flow"}.

    \item[\texttt{module\_dir}] Directory from which the generator module is
    loaded.

    \item[\texttt{module\_class}] Generator filename.  This scenario uses
    \texttt{"BoundaryParticleReservoirPeriodic.py"}.

    \item[\texttt{data\_dir}] Directory receiving generated \texttt{.bin},
    \texttt{.tst}, and \texttt{.rpt} paths.

    \item[\texttt{output\_file\_prefix}] Base output filename.  It also tells
    \texttt{ForceDynamicsBase} which binary file to load when
    \texttt{pdata\_from\_file} is true.

    \item[\texttt{STUDY\_NAME}] Human-readable scenario name used by the
    simulation window and reports.

    \item[\texttt{pdata\_from\_file}] When true, the Python dynamics runner
    loads \texttt{data\_dir/output\_file\_prefix.bin}.  When false, it expects
    inline \texttt{PARTICLE\_DATA} instead.  This reservoir generator normally
    requires true.
\end{description}

\subsection{Particle Preview Plotter}

These fields control the static preview in \texttt{TabFormGenData}.  The axis
limits also control the Python simulation game camera when they are present:

\begin{description}
    \item[\texttt{cell\_range}] Cell-neighborhood range retained by
    \texttt{PlotParticles}.

    \item[\texttt{sphere\_facets}] Number of facets used to draw preview
    spheres.

    \item[\texttt{particle\_range}] Inclusive particle-number range read for
    preview.  A small range can hide boundary particles because mobile
    particles are written first.

    \item[\texttt{particle\_color}] Default preview RGB color.

    \item[\texttt{as\_points}] Draws preview particles as points rather than
    sphere surfaces when true.

    \item[\texttt{vary\_color}] Allows preview color to vary with particle
    type or state when true.

    \item[\texttt{x\_axis\_lims}, \texttt{y\_axis\_lims},
    \texttt{z\_axis\_lims}] Preview axes.  \texttt{SimulationRunner} also
    uses the x and y limits for its game camera, falling back to legacy
    \texttt{Wall*} values when the limits are absent.  The current
    two-dimensional game view does not use \texttt{z\_axis\_lims}.
\end{description}

\subsection{Python Simulation Game UI}

\begin{description}
    \item[\texttt{window\_size}] Pygame window size in pixels, for example
    \texttt{(1000,1000)}.

    \item[\texttt{frame\_rate}] Maximum displayed frames per second.

    \item[\texttt{zoom}] Expands or contracts the game view around the center
    of \texttt{x\_axis\_lims} and \texttt{y\_axis\_lims}.  When those fields
    are absent, the legacy \texttt{WallXMIN/WallXMAX} and
    \texttt{WallYMIN/WallYMAX} rectangle is used.  It does not inspect the
    generated particle distribution.

    \item[\texttt{simulation\_particle\_numbers}] Draws particle numbers in
    the game UI when true.

    \item[\texttt{hsv\_color}] Colors live particles from velocity angle when
    true.

    \item[\texttt{hsv\_sat}] HSV saturation in the interval \([0,1]\).

    \item[\texttt{hsv\_val}] HSV value/brightness in the interval \([0,1]\).

    \item[\texttt{exit\_on\_error}] Closes the simulation on a dynamics error
    when true.  When false, the current runner pauses instead.

    \item[\texttt{run\_debug\_dir}] Required output directory for Python
    simulation reports.

    \item[\texttt{rpt\_frames}] Optional report-frame filter.  If absent,
    reporting uses its normal all-frame behavior.

    \item[\texttt{end\_frame}] Optional automatic final frame.  The caller may
    also supply an end frame directly for batch or headless runs.

\end{description}

\subsection{Dynamics And Lifecycle}

\begin{description}
    \item[\texttt{dt}] Simulation timestep.  It controls integration and is
    also used by the generator to convert required release distance into a
    whole number of frames.

    \item[\texttt{contact\_force\_measure}] Selects \texttt{"depth"} or
    \texttt{"area"}.  The current reservoir configurations select penetration
    depth.  Overlap area remains available for diagnostics and optional force
    comparison.

    \item[\texttt{collision\_stiffness\_q}] Written to each mobile particle's
    \texttt{Data.y}; it scales contact force.  Boundary marker particles are
    currently generated with zero stiffness because they are locality markers,
    not ordinary particle targets.

    \item[\texttt{wall\_contact\_offset}] Configures wall ghost placement as a
    fraction of particle radius.

    \item[\texttt{walls\_on}] Enables wall processing in Python dynamics.

    \item[\texttt{flow\_type}] Enables the \texttt{Data.w} reservoir lifecycle.
    \texttt{"pipe\_reservoir\_entry"} and \texttt{"simple\_reservoir"}
    are recognized.  Other values disable reservoir birth/death gating.

    \item[\texttt{WallXMIN}, \texttt{WallXMAX}, \texttt{WallYMIN},
    \texttt{WallYMAX}] Legacy manual-wall planes.  Manual-wall configurations
    retain them.  Boundary-particle reservoir and CD-nozzle configurations use
    dedicated lifecycle, geometry, and axis-limit fields instead.  The Python
    camera prefers \texttt{x\_axis\_lims} and \texttt{y\_axis\_lims}, with a
    legacy \texttt{Wall*} fallback only for older configurations.

    \item[\texttt{side\_len}] Python shader-flag fallback and cfg value written
    for general cell-array use.  It is distinct from
    \texttt{over\_ride\_side\_length}, which controls this generator's actual
    generated geometry.
\end{description}

\subsection{Reservoir Population And Spacing}

\begin{description}
    \item[\texttt{radius}] Mobile particle radius.

    \item[\texttt{particle\_columns}] Number of release waves generated.

    \item[\texttt{particles\_per\_row}] Number of mobile particles in each
    release wave.

    \item[\texttt{particles\_per\_cell\_row}] Packing target used to derive
    the pipe cross-section in cells.  The generator rejects a configuration
    when this many particle spacings exceeds one cell.

    \item[\texttt{fraction\_of\_diameter\_separation}] Additional
    center-to-center gap as a fraction of diameter.  The generated spacing is:
    \[
    \ell_p=2r(1+f).
    \]

    \item[\texttt{reservoir\_start\_birth}] Birth frame of the first wave.

    \item[\texttt{reservoir\_spacing\_factor}] Required wave clearance in
    particle diameters.  The minimum whole-frame interval is:
    \[
    k=\left\lceil
    \frac{s(2r)}{v_{\mathrm{forward}}\Delta t}
    \right\rceil.
    \]

    \item[\texttt{reservoir\_starting\_vx}] Legacy nominal inlet speed.  It is
    also the default magnitude for \texttt{reservoir\_flow\_speed}.

    \item[\texttt{reservoir\_flow\_speed}] Optional actual velocity magnitude
    used by the angle/alignment velocity generator.  Its default is
    \(\lvert\texttt{reservoir\_starting\_vx}\rvert\).

    \item[\texttt{reservoir\_forward\_speed\_for\_spacing}] Optional
    conservative forward speed used only in the frame-spacing equation.  Its
    default is \texttt{reservoir\_flow\_speed} multiplied by
    \(\max(\texttt{flow\_alignment},0.1)\).

    \item[\texttt{reservoir\_birth\_jitter\_frames}] Adds random per-particle
    birth jitter from zero through this value.  The current generator also adds
    this maximum to the base wave step, so consecutive wave bases are separated
    by \(k+\texttt{jitter}\).

    \item[\texttt{reservoir\_random\_seed}] Makes position, direction, and birth
    jitter reproducible.

    \item[\texttt{reservoir\_position\_jitter\_fraction}] Fraction of the
    available inter-particle gap used for random cross-axis position jitter.
    It does not jitter by a fraction of the full particle diameter.

    \item[\texttt{flow\_angle}] Preferred velocity direction in degrees.
    Zero degrees points along positive \(x\).  This is independent of wall
    orientation.

    \item[\texttt{flow\_alignment}] Blend between preferred direction and a
    random unit direction.  One is fully aligned; zero is fully random.  The
    blended direction is normalized before multiplying by flow speed.

    \item[\texttt{reservoir\_random\_lateral\_velocity}] Selects random legacy
    lateral velocity when true.

    \item[\texttt{reservoir\_lateral\_velocity}] Magnitude used by the legacy
    lateral-velocity path.  When \texttt{flow\_alignment < 1}, the angle-based
    generated lateral component replaces this legacy value.

    \item[\texttt{reservoir\_inlet\_x}, \texttt{reservoir\_inlet\_y}]
    Dedicated runtime birth planes.  ForceDynamics uses the field for the
    selected flow axis and falls back directly to the corresponding legacy
    \texttt{Wall*} minimum or maximum when it is absent.  Generic
    \texttt{inlet\_x}/\texttt{inlet\_y} aliases are not used.

    \item[\texttt{reservoir\_outlet\_x}, \texttt{reservoir\_outlet\_y}]
    Dedicated center-crossing retirement planes.  ForceDynamics falls back
    directly to the corresponding legacy \texttt{Wall*} value when the field
    for the selected flow axis is absent.  Generic
    \texttt{outlet\_x}/\texttt{outlet\_y} aliases are not used.

    \item[\texttt{reservoir\_base\_x}, \texttt{reservoir\_base\_y}]
    Optional cross-axis row origins before the generator adds its centering
    offset.  Defaults are \(1.0\).

    \item[\texttt{reservoir\_z}] Optional particle plane coordinate.  The
    default is \(2.0\).
\end{description}

\subsection{Boundary-Particle Geometry}

\begin{description}
    \item[\texttt{boundary\_particle\_function}] Selects generator orientation
    and the wall evaluator.  \texttt{"horizontal\_wall"} creates horizontal
    pipe flow with upper/lower marker rows.  \texttt{"vertical\_wall"}
    creates vertical flow with left/right marker columns.

    \item[\texttt{wall\_type}] Written to the generated \texttt{.tst} file for
    the Vulkan wrapper.  Python \texttt{ForceDynamics} currently recognizes
    boundary-particle walls by finding particles with \texttt{ptype=1}, not by
    branching on this string.

    \item[\texttt{over\_ride\_side\_length}] If positive, directly selects the
    generator's cubic cell-array side length.  If absent or nonpositive, the
    generator uses \texttt{side\_len}.  Derivation from row packing plus a
    hardcoded three-cell guard on each side is retained only for legacy cfg
    files that provide neither positive field.

    \item[\texttt{wall\_boundary\_length}] Explicit marker-wall length
    along the flow axis.  It is required and must be positive for
    \texttt{BoundaryParticleReservoirPeriodic}; marker geometry is no longer
    derived from \texttt{WallXMIN...WallYMAX}.

\end{description}

The boundary marker radius is currently hardcoded to \(0.25\), marker
\texttt{ptype} is \(1\), marker \texttt{Data.w} is zero, and marker stiffness
is zero.  Each marker writes its evaluator ID through \texttt{pdata.temp\_vel}
to runtime \texttt{Data.z}: horizontal is 1, vertical is 2, and CD nozzle is
3.  These are current implementation facts, not configurable fields.

\subsection{Generated Test Metadata}

The generator calculates \texttt{dispatchx} as the null particle plus the
number of mobile particles; it is not a cfg input for this scenario.  Boundary
markers are stored after mobile particles and are included in
\texttt{num\_particles}, but they are not compute sources.

The following cfg fields are copied into the generated \texttt{.tst} file:

\begin{verbatim}
dispatchy, dispatchz
workGroupsx, workGroupsy, workGroupsz
cell_occupancy_list_size
\end{verbatim}

The boundary-particle reservoir generator no longer writes the retired
zero-valued collision-density compatibility fields.

\subsection{Generator Simulation Checks}

Both boundary-particle generators run a preflight check after calculating
their geometry and before writing particles.  Manual configuration values are
compared with the calculated simulation domain.  Generation stops with one
aggregated error when a physical inconsistency is found.

For \texttt{BoundaryParticleReservoirPeriodic}, the checks include positive
particle and timestep values, valid population counts, cell-row packing,
required inlet/outlet fields, inlet agreement with the calculated inlet marker
wall, an outlet beyond the terminal marker wall, and all lifecycle planes
inside the cell domain.  The birth center must clear the inlet wall by at least
the configured minimum diameter factor.

For \texttt{BoundaryCDNozzleReservoir}, the checks additionally require valid
dimension, positive section lengths and radii, a throat radius no larger than
the inlet or exit radius, and:
\[
\texttt{wall\_boundary\_length}=
L_{\mathrm{inlet}}+L_{\mathrm{converge}}+L_{\mathrm{throat}}+
L_{\mathrm{diverge}}+L_{\mathrm{exit}}.
\]
The lifecycle outlet must lie beyond the calculated nozzle end while remaining
inside the cell domain.  Axis limits that exclude valid generated geometry are
reported as warnings because they affect visibility rather than dynamics.

\subsection{Important Cross-System Effects}

The boundary-particle cfg families now declare dedicated migration fields:

\begin{verbatim}
reservoir_inlet_x, reservoir_outlet_x
wall_boundary_length
nozzle_start_x, nozzle_center_y  // CD nozzle only
\end{verbatim}

The boundary-particle reservoir and CD nozzle cfg families no longer define
legacy \texttt{WallXMIN...WallYMAX}.  Manual-wall cfg families retain those
fields unchanged.  Display bounds are shared through
\texttt{x\_axis\_lims} and \texttt{y\_axis\_lims} rather than duplicated as
separate game-view fields.

\begin{enumerate}
    \item The include file overrides the parent cfg.  Always inspect the merged
    values before diagnosing a run.

    \item \texttt{over\_ride\_side\_length} can center generated particles far
    from the original cfg wall rectangle.  Both preview and game UI use the
    configured axis limits, so those limits must cover the generated geometry
    or valid particles can remain offscreen.

    \item \texttt{side\_len} and \texttt{over\_ride\_side\_length} are not
    aliases.  Conflicting values can make Python runtime bounds, generated test
    metadata, and particle geometry describe different domains.

    \item \texttt{flow\_angle} controls velocity direction; it does not select
    wall orientation.  \texttt{boundary\_particle\_function} selects the
    current straight-wall orientation.

    \item Setting \texttt{display\_base="live"} does not disable shadow
    computation.  Set \texttt{shadow\_base="none"} in
    \texttt{ParticleUtil.cfg} when only one model should run.
\end{enumerate}
